% 本文件由 LaTeX 排版 Agent 根据有效 translation.json 自动生成。
% author=Tertullian; work_id=0308; title=答犹太人

\chapter{答犹太人}

% structure: p0001 source=h1 target=omitted reason=page-title-already-rendered-by-parent

% structure: p0002 source=h2 target=section reason=relative-heading-rank; base=h2

\section{第一章 写作缘由。犹太人与外邦人的相对地位示例。}

最近发生了一场辩论，在一位基督徒与一位犹太皈依者之间展开。他们以好辩的绳索交替拉扯，将白日拖延至黄昏。此外，由于某些支持者的对立喧嚷，真理开始被一层迷雾遮蔽。因此，我们乐意将那因争论的混乱喧嚣而无法逐点充分阐明之事，加以更仔细的研读，并以笔墨来判定那些已付诸讨论的问题，以供阅读。\par

事实上，为外邦人求得神圣恩宠的机缘，因着那自命维护天主的律法的人并非犹太人，而是外邦人，并非\Quote{以色列子孙}的犹太人，而获得了卓越的适宜性。\BibleRef{斐 3:5} 因为这一事实——外邦人得以进入天主的律法——已足以阻止以色列人自夸，以为\Quote{外邦人被算为桶里的一滴水}，或如\Quote{打谷场上的尘土}；尽管我们有天主亲自作为充足的担保者和信实的应许者，因祂曾应许亚巴郎说：\Quote{地上万民都要因你的后裔蒙福}；并且从黎贝加的胎中\Quote{将要生出两个民族和两个邦国}——当然是指犹太人的，即以色列；以及外邦人的，即我们的。两者皆被称为\Emph{民族}和\Emph{邦国}；以免有人从这名义的称呼妄自为自身索取恩宠的特权。因为天主定下\Quote{两个民族和两个邦国}将从一个妇人的胎中生出：恩宠在这名义的称呼上并未加以区分，而是在出生的次序上；使得那先出胎的，要服从于\Quote{较小的}，即后出的。因为天主如此对黎贝加说：\Quote{两个民族在你腹中，两个邦国要从你怀中分出；这民族要胜过那民族，年长的要服事年幼的。} 既然犹太人的\Emph{民族}或\Emph{邦国}在时间上居先，并因律法中首蒙宠爱的恩宠而\Quote{较大}，而我们则在时代的年龄中被视为\Quote{较小}，因为我们在世界末世才获得神圣慈悲的认识；那么，无疑地，按照天主圣言的谕令，那\Emph{在先}而\Quote{较大的}民族——即犹太民族——必定要服事那\Quote{较小的}；而那\Quote{较小的}民族——即基督徒——要胜过那\Quote{较大的}。因为，按照神圣经卷的记载，犹太人的\Emph{民族}——即更古老的——完全离弃了天主，卑贱地事奉偶像，离弃了神性，降服于神像；而\Quote{百姓}对亚郎说：\Quote{为我们造神祇，在我们前面引路。} 当妇女的项链和男子的指环中的金子完全被火熔化，并现出一个牛犊般的头时，以色列人一致（离弃天主）尊崇这偶像，说：\Quote{这些神祇是将我们从埃及地领出来的。} 因为后来在诸王治理他们的时代，他们又与雅洛贝罕一同崇拜金犊、木偶，并奴役自己于巴耳。由此证明，他们从神圣经卷的篇章中常被描绘为犯有偶像崇拜的罪行；而我们\Quote{较小的}——即后起的——\Emph{民族}却离弃了昔日奴役地事奉的偶像，归向了同一个天主，而那以色列人如上所述已离弃了祂。\BibleRef{得前 1:9} 因为正是如此，那\Quote{较小的}——即后起的——\Emph{民族}胜过了那\Quote{较大的民族}，因它获得了神圣宠爱的恩宠，而以色列却与此恩宠隔绝了。\par

% structure: p0005 source=h2 target=section reason=relative-heading-rank; base=h2

\section{第二章 律法先于梅瑟。}

因此，让我们面对面站定，在确定的界限内判定这实际问题的实质。\par

因为为何要相信，宇宙的建立者、整个世界的治理者、人类的塑造者、万国的播种者天主，透过梅瑟只给予一个民族法律，而不说祂将它分派给所有民族？因为除非祂已将它给予所有人，否则祂绝不会惯常允许甚至外邦中的归化者接近它。但是——与天主的良善和祂的公义相称，作为人类的塑造者——祂将同一法律赐给了所有民族，祂曾命令在确定和规定的时间内遵守它，当祂愿意时，透过祂愿意的人，并按祂愿意的方式。因为在世界之初，祂就给了亚当和厄娃本身一条法律，即他们不可吃乐园中央那棵树上的果子；但如果他们反其道而行，他们就要因死亡而死亡。假如这条法律被遵守，它本已足够他们。因为在这条赐给亚当的法律中，我们认出了所有后来透过梅瑟赐下的诫命的雏形；即：\Quote{你应全心、全灵爱上主你的天主；你应爱人如己；不可杀人；不可奸淫；不可偷盗；不可作假见证；应孝敬父母；以及，不可贪恋别人的财物。}因为原始的法律是在乐园中赐给亚当和厄娃的，作为所有天主诫命的母胎。简言之，如果他们爱上了主、他们的天主，他们就不会违反祂的诫命；如果他们惯常爱了近人——即他们自己——他们就不会相信蛇的劝说，因而不会因违犯天主的诫命、从永生中堕落而自杀；他们也会避免偷盗，如果他们不是偷吃树上的果子，也不会急于躲在一棵树下来逃避主、他们的天主的目光；他们也不会与发假誓的魔鬼为伍，相信他所说的他们要\Quote{如同天主}；从而他们也不会得罪天主，作为他们的父亲，祂曾用地上的泥土塑造了他们，如同从母胎中；如果他们不贪恋别人的东西，他们就不会吃那违禁的果子。\par

因此，在这普遍和原始的天主法律中，祂曾以树果为例，制定了其遵守，我们认得所有后来法律的诫命都\Emph{内}含在其中，当它们在适当的时候被\Emph{展}现出来时。因为后来法律之附加，是那先前已预设了诫命的同一位的工作；因为这同样是祂的权限，祂先前已决意要塑造公义的受造物，后来再加以训练。因为，如果祂创始了纪律，祂又扩展纪律，有什么可怪？如果祂开始了，祂又推进，有什么可怪？总之，在写在石版上的梅瑟法律之前，我主张有一部不成文的法律，它被自然地理解，并被先祖们惯常遵守。因为诺厄是凭什么被\Quote{找到为义人}，如果在他身上自然法律之义德没有先行？亚巴郎凭什么被算作\Quote{天主的朋友}，如果不是基于（遵守）自然法律的公平和义德？默基瑟德凭什么被称\Quote{至高天主的司祭}，如果，在肋未法律司祭职之前，没有肋未人惯常向天主献祭？因为这样，在上述圣祖之后，法律在（众所周知的）他们出离埃及之后，经过四百年的间隔和空间，赐给了梅瑟。事实上，正是在亚巴郎的\Quote{四百三十年}之后，法律才被赐下。由此我们明白，天主的法律在梅瑟之前就已存在，并非首先（赐下）在曷勒布，也不是在西乃和旷野，而是更古老；（首先）存在于乐园中，随后为圣祖们作了修订，再为犹太人，在确定的时期；因此，我们不应将梅瑟的法律视为原始的法律，而应视为后续的法律，是天主在确定的时期也向外邦人颁布，并且，在多次透过先知应许之后，又加以改善；并且预先警告说，将发生这样的事，正如\Quote{法律是藉着梅瑟而赐予的}\BibleRef{若 1:17}，在确定的时间，同样应相信它是暂时被遵守和持守的。并且，我们不要废除天主所具有的这种权能，它根据时代的情况调整法律的诫命，以促进人的救恩。总之，让那些主张安息日仍应作为救恩的药膏而被遵守，以及因死亡的威胁而在第八日行割损的人，教导我们，在过去的时间，义人是否遵守了安息日或行了割损，并因此成了\Quote{天主的朋友}。因为如果割损洁净人，而天主造了未受割损的亚当，为什么祂没有在他犯罪后给他行割损，如果割损能洁净呢？无论如何，在安置他在乐园时，祂指派了一个未受割损的人作乐园的居民。因此，既然天主创造了未受割损和不守安息日的亚当，他的后代亚伯尔，未受割损和不守安息日，向他献祭，就得到了他的赞许；因为他以纯朴的心接受了他所献的，并且拒绝了其弟加音的祭品，因为他没有正确地剖分他所献的。诺厄，也是未受割损的——是的，也不守安息日——天主却把他从洪水中救出。哈诺客，也是极正义的人，未受割损和不守安息日，天主将他从这个世界移去；他没有先尝到死亡，为的是，作为永生的候选人，他此时可以指示我们，我们也能没有梅瑟法律的负担而中悦天主。默基瑟德，\Quote{至高天主的司祭}，未受割损和不守安息日，却被选为天主的司祭。罗特，亚巴郎的兄弟，也证明是由于义德的功劳，没有遵守法律，他才从索多玛人的火灾中被救出。\par

% structure: p0009 source=h2 target=section reason=relative-heading-rank; base=h2

\section{第三章 论割礼与旧法律的废止。}

但是，你或许会说，亚巴郎受过割损。不错，他在受割损前就已中悦了天主，而且他也没有遵守安息日。因为他领受了割损，但那只是为那时期作\Quote{标记}，并非获得救恩的特权。事实上，后来的族长如默基瑟德就没有受过割损，他未受割损，却向已受割损的亚巴郎在作战归来时献上饼和酒。\newline{}\Quote{可是，}你或许会说，\Quote{梅瑟的儿子曾一度差点被天使扼死，若不是漆颇辣用石子割去婴孩的包皮的话；因此，\InnerQuote{若有谁未割去他肉身的包皮，就有极大的危险。}}不！如果割损果真带来救恩，梅瑟本人就不会对自己的儿子忽略在第八日行割损；而众所周知，漆颇辣是在旅途中因天使的催逼才行割损的。因此，我们要考虑，一个婴孩被迫的割损并不能为所有民族立下规条，并建立一条遵守此诫命的法律。因为天主预见到祂将赐予以色列民这种割损作为\Quote{标记}，而非为救恩，就因此催逼梅瑟的儿子——他们未来的领袖——行割损；因为祂通过梅瑟开始将割损的诫命赐给百姓，百姓不应因看到在他们领袖的儿子身上已显著出现的（疏忽）榜样而轻视它。因为割损必须被赐予，但只是作为\Quote{标记}，以便以色列在末世时能被区分出来，那时根据他们的所作所为，他们将被禁止进入圣城，正如我们通过先知的话所看到的：\BibleQuote{你的土地荒凉；你的城邑被火焚毁；你的田地在你眼前被外邦人吞吃；被外邦人倾覆荒废，熙雍女子被撇下，像葡萄园里的草棚，像黄瓜田中的茅屋，又像被攻打的城邑。}为何如此？因为先知后续的话责备他们说：\BibleQuote{我养育了儿女，使他们长进，他们却背弃了我；}又说：\BibleQuote{你们若伸出双手，我必掩面不看你们；你们若多多祈祷，我也不听：因为你们的手沾满了血；}\BibleRef{依 1:15}又说：\BibleQuote{祸哉！犯罪的民族，充满罪恶的百姓，邪恶的儿子！你们完全离弃了天主，惹起了以色列的圣者的忿怒。}\BibleRef{依 1:4}因此，这是天主的远见——赐予以色列割损作为标记，以便在时候来到时，他们能因上述的罪行而被禁止进入耶路撒冷；这情况因为将要发生而曾被宣告；又因为我们看到它已实现，就被我们认出来。因为肉身的割损是暂时的，曾在一个顽梗的百姓中被作\Quote{标记}，而心灵的割损则赐给一个顺从的百姓以得救恩；正如先知耶肋米亚所说：\BibleQuote{你们当为自己开垦田地，不要撒种在荆棘中；当受割损归于天主，割去你们心中的包皮；}在另一处又说：\BibleQuote{看，时日将到，上主说，我必与犹大家和以色列家订立新约，不像我昔日握住他们的手，领他们出离埃及时，与他们的祖先所订立的约。}由此我们明白，从前所赐的割损将要停止，以及新法律（不像祂已赐给祖先的那样）的来临，都被宣告了；一如依撒意亚预言说，在末日上主的殿宇必要矗立在群山之上：\BibleQuote{它必被高举，}他说，\BibleQuote{超越山丘；万民都要涌向它；许多人要去，且说：来，我们登上上主的山，到雅各伯的天主的殿里，}\BibleRef{依 2:2}——不是\Emph{厄撒乌}，那先前的儿子，而是\Emph{雅各伯}，那后来的；也就是说，\Emph{我们}这\Quote{民族}的，它的山是基督，\BibleQuote{非人手所凿的石头，充满全地}，见于达尼尔书。简言之，这新的法律从\BibleQuote{雅各伯的天主的殿}出来，依撒意亚在接下来的话中宣告说：\BibleQuote{因为法律将出自熙雍，上主的话将出自耶路撒冷，且要统治万民，}——即我们，从万民中被召叫的——\BibleQuote{他们将把刀剑铸成犁头，把枪矛制成镰刀；民族不再举刀相攻，也不再学习战事。}\BibleRef{依 2:3}此外，还有谁被理解为就是我们呢？我们受新法律的充分教导，遵守这些行为——旧法律被废除，其废除的来临由行为本身证明？因为旧法律的惯常用刀剑复仇，以\BibleQuote{眼还眼}，对伤害施以报复。但新法律的惯常是指向仁慈，将\Quote{刀剑}和\Quote{枪矛}的原始凶暴转变为和平，并将原先对法律的敌人和对手\Quote{作战}的行为改造为犁地和耕种的和平行为。因此，正如我们上面所表明的，旧法律和肉身的割损即将停止，同样，新法律的遵守和心灵的割损已在自愿的和平顺从中照耀出来。因为他说：\BibleQuote{我不认识的民族事奉了我；他们一听到就服从了我。}先知们宣告了这事。但那个不认识天主的\Quote{民族}是谁呢？不就是我们吗？我们在往昔不认识天主。谁在耳闻中听从了祂？不就是我们吗？我们离弃了偶像，归向了天主。因为以色列——曾为天主所认识，曾被祂在埃及\Quote{抚养}，渡过红海，在旷野中四十年以玛纳养育，被塑造成永恒的样式，没有沾染人的情欲，没有吃这世界的食物，而是吃\BibleQuote{天使的粮食}——玛纳——并因天主的恩惠而紧紧与祂相连——却忘记了祂的上主和天主，对亚郎说：\BibleQuote{为我们造神像，在我们前面引路；因为那个领我们出埃及地的梅瑟，不知道遭遇了什么事。}因此，我们这些\BibleQuote{从前不是天主的子民}的，如今成了祂的子民，因接受了上述的新法律和预言过的新割损。\par

% structure: p0011 source=h2 target=section reason=relative-heading-rank; base=h2

\section{第四章 论守安息日。}

因此，既然肉身的割损和旧律的废除已在特定时期完成，那么安息日的遵守也同样是暂时的。\par

因为犹太人声称，天主从起初就圣化了第七天，因为祂在那天歇了祂所做的一切工作；并且因此，梅瑟也对人民说：\Quote{应记住安息日，守为圣日：不可做任何劳役，惟有关乎性命的事除外。} 由此我们（基督徒）明白，我们更应常常遵守远离一切 \Quote{劳役} 的安息，而不只每第七天，乃是通过一切时间。由此为我们产生了问题：天主愿意我们遵守的是什么样的安息？因为圣经指示了一个永久的安息和一个暂时的安息。因为依撒意亚先知说：\Quote{\Emph{你们的}安息日，我的灵魂憎恨；} \BibleRef{依 1:13} 在另一处又说：\Quote{\Emph{我的}安息日，你们亵渎了。} 由此我们看出，暂时的安息是属于人的，永久的安息被认为是天主的；关于这安息，祂通过依撒意亚预言说：\Quote{且要有，} 祂说，\Quote{月复一月，日复一日，安息复安息；凡有血肉的都要来耶路撒冷朝拜，这是主说的；} 我们理解，这已在基督的时代应验了，那时 \Quote{凡有血肉的}——即万民——\Quote{都来耶路撒冷朝拜} 天主圣父，通过祂的儿子耶稣基督，正如先知所预言的：\Quote{看，通过我，归依者要到祢那里去。} 因此，在这个暂时的安息之前，已有一个永久的安息被预示和预言；正如在肉身的割损之前，已有一个精神的割损被预示。简言之，让他们教导我们，正如我们已经预先说明的：亚当是否遵守了安息日；或亚伯尔，在向天主献上圣洁的祭品时，是否因对安息日的宗教敬畏而中悦了祂；或哈诺客，在被提时，是否是安息日的遵守者；或建造方舟的诺亚，是否因洪水而守了一个漫长的安息日；或亚巴郎，是否在遵守安息日时献了依撒格他的儿子；或默基瑟德，在他的司铎职中，是否接受了安息日的律法。\par

但犹太人必定会说，自从这诫命通过梅瑟赐下以来，遵守它就是有约束力的。因此显而易见，这诫命不是永久的，也不是精神的，而是暂时的，有一天会停止。简言之，这庆典不应当在于安息日——即第七天——不做工，这是如此真实，以至于农的儿子若苏厄，在攻打耶里哥城时，声称他从天主领受了命令：命令人民，司铎要抬着天主的约柜，绕城七天；这样，当第七天的绕行完成后，城墙会自行倒塌。\BibleRef{苏 6:1} 这事就这样成了；当第七天的空间结束时，正如所预言的，城墙倒塌了。由此明显可见，在七天的数目中，有一个安息日。因为七天，无论从何时开始，必然包含一个安息日；在那一天，不仅司铎必须工作，而且全以色列人民要用刀剑掠夺那城。毫无疑问，他们 \Quote{做了劳役}，因为遵从天主的命令，他们夺取了战利品。因为在玛加伯时代，他们也在安息日勇敢作战，击溃外敌，并通过在安息日作战，将祖传的法律恢复到原始的生活方式。我不认为他们如此辩护的是别的律法，而是他们记得其中存在关于 \Quote{安息日} 的诫命。\par

由此可见，此类诫命的力量是暂时的，与当前情况的需要有关；天主从前赐给他们这样的律法，并不是为了永久的遵守。\par

% structure: p0016 source=h2 target=section reason=relative-heading-rank; base=h2

\section{第五章 论祭献。}

同样，我们又证明，用地上祭品所作的祭献和属神的祭献曾被预言；而且，从起初，地上的祭献就已在加音身上预表为\Quote{长子}的祭献，即以色列的祭献；而相反的祭献则被证明是\Quote{幼子}亚伯尔的祭献，即\Emph{我们}民族的祭献。因为长子加音从地里的出产奉献礼物给天主；但幼子亚伯尔从他羊群的首生奉献。\BibleQuote{天主看中了亚伯尔和他的礼物；却没有看中加音和他的礼物。天主对加音说：你为什么垂头丧气？你若的确献得对，却没有分得对，难道不犯罪吗？你当安静。因为你的归依必归于你，他必辖管你。然后加音对亚伯尔说：我们到田间去；他就与他在那里，把他杀了。天主对加音说：你的兄弟亚伯尔在哪里？他说：我不知道；我岂是看守我兄弟的吗？天主对他说：你兄弟的血的声音从地里向我哀告。因此，地受了诅咒，它开口从你手里接受了你的兄弟的血。你必在地上流离失所，凡遇见你的必杀你。}从这个事件中我们得知，\Quote{众民}的双重祭献从起初就已预表。简言之，当司祭的法律通过梅瑟在\BibleBook{肋未}{纪}中制定时，我们发现给以色列人规定：祭献不可在任何其他地方献于天主，只可在应许之地；那是主天主将要赐给以色列民和他们的弟兄的，以便以色列人进入那里后，在那里举行祭献和全燔祭，为罪也为灵魂；且只在圣地，别无他处。那么，为什么圣神后来通过先知预言，将来在每一个地方和每一块土地上都要向天主献祭呢？正如祂通过十二先知之一的玛拉基亚天使所说：\BibleQuote{我不会从你们手中接受祭献；因为从日出之地到日落之地，我的名已在万民中闻名，全能的上主说：并且在各处，他们向我名奉献洁净的祭献。}再者，达味在\BibleBook{}{圣咏}中说：\BibleQuote{列邦的国土啊，你们当献给天主}——无疑是因为宗徒的宣讲必须\Quote{传到}\Quote{每一块土地}——\BibleQuote{将荣耀和尊贵献给天主；将祂名的祭献献给天主：举起牺牲，进入祂的庭院。}因为向天主奉献不是藉着肉身的祭献，而是藉着属神的祭献，我们这样读到，正如经上所记：痛悔和谦卑的心是天主所悦纳的牺牲；又在别处说：\BibleQuote{向天主献上赞美的祭，并向至高者还你的愿。}因此，属神的\Quote{赞美的祭献}被指出，而\Quote{痛悔的心}被证明是天主所悦纳的祭献。同样，正如肉身的祭献被理解为被摒弃——依撒意亚也如此说：\BibleQuote{你们众多祭献对我何益？上主说}\BibleRef{依 1:11}——这样，属神的祭献被预言为蒙悦纳，正如先知所宣告。因为，\BibleQuote{即使你们给我带来}祂说，\BibleQuote{上好的细面，也是徒然的恳求之礼：是我所憎恶的；}祂又说：\BibleQuote{你们的全燔祭和祭献，以及公羊的脂油和公牛的血，我都不想要，即使你们来让我看见，我也不要：因为谁从你们手中要求这些东西呢？}因为\BibleQuote{从日出之地到日落之地，我的名在万民中已经闻名，上主说。}但对于属神的祭献，祂补充说：\BibleQuote{并且在各处，他们向我名奉献洁净的祭献，上主说。}\par

% structure: p0018 source=h2 target=section reason=relative-heading-rank; base=h2

\section{第六章 论旧律法的废除与废除者}

因此，既然明显可见：一个暂时的安息日曾显示，一个永久的安息日已被预告；一个肉身的割礼曾被预告，一个属神的割礼已被预先指明；一个暂时的律法和一个永久的律法被正式宣告；肉身的祭献和属神的祭献曾被预表；那么，这一切诫命在先前时期以肉身方式给了以色列人之后，必然要有一个时期来到，使古律法和旧礼仪的诫命终止，新律法的应许、属神祭献的承认以及新约的应许将接续而来；同时，从高处来的光将照亮我们这些坐在黑暗中、被死亡的阴影拘留的人。因此，我们负有一种必要\BibleRef{格前 9:16}，既然我们已先提出新律法曾被先知预言，且不是像当初祂领他们出埃及时已赐给他们祖先的那种律法，我们就必须表明并证明：一方面，那旧律法已经终止；另一方面，所应许的新律法如今正在运行。\par

事实上，我们首先必须探究：是否有一位新法的赐予者、新约的继承者、新祭献的司铎、新割损的洁净者、永恒安息日的遵守者，祂要废除旧法，建立新约，奉献新祭献，废止古礼，连同旧割损和它的安息日一起废除，并宣告那不朽的新王国。我说，我们必须探究：这位新法的赐予者、属神安息日的遵守者、永恒祭献的司铎、永恒王国的永恒统治者，是否已经来临？如果祂已来临，我们就当事奉祂；如果祂尚未来临，我们就当等候祂，直到祂的来临显明旧法的诫命已被废除，新法的开端应当兴起。首先，我们必须确认：古法和先知不可能终止，除非那通过同一法律和同一先知不断被宣告要来临的那位已经来临。\par

% structure: p0021 source=h2 target=section reason=relative-heading-rank; base=h2

\section{第七章 论基督是否已经降临的问题}

因此，我们在这个问题上针锋相对：那不断被宣告要来临的基督是否已经来临，还是祂的来临仍是希望的对象？为证明这个问题本身，我们也必须查考先知们宣告基督将要来临的时期；如果我们能确认祂是在那些时期内来临的，我们就可以毫无疑问地相信，祂正是那位先知之歌永远颂扬的未来来临者，我们——即外邦人——曾注定要相信祂；而当我们一致承认祂已来临，我们无疑也会相信新法已由祂赐下，并承认在祂内并藉着祂为我们订立的新约。因为我们知道，连犹太人也无法否认基督要来，因为他们将自己的希望寄托于祂的来临。关于此事，我们无需更多探究，因为古时所有的先知都已预言过；如依撒意亚说：\Quote{主天主如此对我的主基督说：我紧握祂的右手，使万民听从祂；我要击碎君王的力量；在祂面前开启城门，城门不再关闭。}我们正看到这事应验。天主圣父紧握谁的右手，不就是祂的儿子基督的吗？万民都听从了祂，即万民都相信了祂——此外，祂的宣讲者宗徒们，在达味的圣咏集中被指出：\Quote{他们的声音传遍普世，}他说，\Quote{他们的话传到地极。}因为还有谁使万民相信，除了那已来临的基督？万民相信了谁？帕提亚人、玛待人、厄蓝人、住在美索不达米亚的、亚美尼亚、夫黎基雅、卡帕多细雅，以及本都、亚细亚、潘菲利亚的居民，寄居埃及的、住在基勒乃以外的阿非利加地区的、罗马人和侨民，还有耶路撒冷的犹太人，以及所有其他民族；例如，时至今日，革突利人的不同种族、摩利人的多样区域、西班牙的所有边界、高卢的各民族、不列颠人的居所——对罗马人难以进入，却屈服于基督；还有撒马提人、达西亚人、日耳曼人、西徐亚人，以及许多遥远的民族、许多我们不知且几乎无法一一列举的省份和岛屿。在所有这些地方，已来临的基督之名作王，正如在祂面前，所有城门都已敞开，无门关闭，在祂面前铁闩粉碎、铜门敞开。虽然这些表达也有属灵的含义——即被魔鬼以各种方式封锁的个人之心，被基督的信德打开——但它们在实际上显然已经应验，因为在所有这些地方都住着基督之名的\Quote{子民}。因为除了基督、天主子，谁能统治万民？祂曾被宣告要永远统治万民。如果撒罗满\Quote{统治}，那只是在犹太境内：从贝尔舍巴直到丹，是祂王国的边界。如果达理阿统治巴比伦人和帕提亚人，他没有统治万民；如果法郎或任何继承他在埃及世袭王国的人，他只在埃及境内拥有王权；如果拿步高和他的小王们，从印度到埃塞俄比亚是他王国的边界；如果亚历山大马其顿人，他在完全征服之后，也不过统治了整个亚细亚和其他地区；如果日耳曼人，至今他们仍不得越过自己的边界；不列颠人被自己的海洋包围；摩利人的民族和革突利的野蛮人被罗马人封锁，以免超出他们自己的区域。关于罗马人自己，我又能说什么呢？他们用自己军团驻防巩固帝国，却无法将王权扩展到这些民族之外。但基督的名字正在扩展，被到处信仰，被上述所有民族崇拜，到处为王，到处受敬拜，平等地赐予所有人。没有哪个国王在祂面前更受偏爱，没有哪个野蛮人更少喜乐；没有尊贵或家世享有功绩的区别；对所有人祂是平等的，对所有人是王，对所有人是审判者，对所有人是\BibleQuote{天主和主}\BibleRef{若 20:28}。你看到这事正在发生，就不会怀疑我们所断言的。\par

% structure: p0023 source=h2 target=section reason=relative-heading-rank; base=h2

\section{第八章 论基督诞生、受难及耶路撒冷毁灭的时期}

因此，必须探究所预言和未来的基督的诞生、祂的受难以及耶路撒冷城被消灭（即其荒废）的时期。因为达尼尔说：\Quote{圣城和圣地都要与来临的首领一同被消灭，顶楼被毁坏到底。} 因此，必须探究来临的首领基督的时期\BibleRef{依 55:4}，我们将在达尼尔书中追溯这些时期；在计算之后，我们将证明祂已经来临，甚至依据规定的时期以及祂的适当征兆和作为。而基于祂来临后所宣布的后果，我们再次证明这些事，为使我们相信一切既已应验如所预见。\par

因此，达尼尔如此预言了祂，既显示祂将在何时解放万民，也显示在基督受难后，那座城如何要被消灭。因为他如此说：\Quote{在玛待裔的阿哈随鲁的儿子达理阿元年，他统治了加色丁人的王国，我达尼尔在经卷中明白了年数……当我还在祈祷中说话时，看，那人加俾额尔，我在起初的异象中见过他，飞着，在晚祭的时刻触摸了我，使我明白，并对我说话，说：达尼尔，我现在出来使你领悟；在你的恳求开始时就发出了一句话。我来向你宣布，因为你是个蒙受爱慕的人；且要思量这话，明白这异象。七十个星期已缩减在你的人民和圣城上，直到过犯成熟，罪恶被封，义德因恳求而获得，永久的义德被引入；并且为封住异象和先知，并傅油至圣的圣者。你要知道、彻底看见并明白：从发出重建耶路撒冷的话直到首领默西亚，是七个半星期和六十二个半星期；它将转变，被建造成高台和壕沟，时期将被更新；在这六十二个半星期之后，傅油者将被消灭，不复存在；城和圣地将与来临的首领一同被消灭；他们像洪水一样被剪除，直到战争的终结，这战争将被剪除至毁灭。他将在许多人中坚定盟约。在一个星期和半个星期中，我的祭品和奠品将被取走，在圣地中，荒凉的可憎之物，直到终结的时期，关于这荒凉将被给予。}\par

因此，我们注意界限——他事实上预言了七十个星期，\Emph{在这期间}若他们接待祂，\Quote{它将被建造成高台和壕沟，时期将被更新。} 但是天主预见将要发生的事——他们不仅不接待祂，反而将迫害并交付祂至死——就总括了，并说，在六十二个半星期中祂诞生，至圣的圣者被傅油；但在七个半星期完成时，祂必须受难，圣城在一个半星期后要被消灭——也就是七个半星期已经完成。因为他如此说：\Quote{城与圣地要与来临的首领一同被消灭；他们像洪水一样被剪除；他必毁坏顶楼到底。} 因此，我们从何处证明基督在六十二个半星期内来临了呢？此外，我们从达理阿元年计数，因为在这特定时刻向达尼尔显示了这特定异象；他说：\Quote{要明白并推测，在完成你的话语时，我给你这些答复。} 因此我们必须从达理阿元年（达尼尔见这异象时）开始计算。\par

因此，让我们看年数如何被填满直至基督来临：\par

因为达理阿在位……十九年。\newline{}阿尔塔薛西斯在位……四十一年。\newline{}然后国王奥古斯（也称为居鲁士）在位……二十四年。\newline{}亚尔古……一年。\newline{}另一位达理阿，也称为美拉斯……二十一年。\newline{}马其顿的亚历山大……十二年\par

然后，在亚历山大之后——他曾统治玛待人和波斯人，征服了他们，并在亚历山大里亚稳固地建立了他的王国，他甚至以自己的名字命名那城——在他之后，（在亚历山大里亚）统治了：\par

索特尔……三十五年。\newline{}继他的是斐拉德尔福斯，在位……三十八年。\newline{}继他的是厄厄尔革忒斯……二十五年。\newline{}然后是斐罗帕托……十七年。\newline{}他之后是厄丕法乃……二十四年。\newline{}然后是另一位厄厄尔革忒斯……二十九年。\newline{}然后是另一位索特尔……三十八年。\newline{}托勒密……三十七年。\newline{}克娄巴特拉……二十年零五个月。\newline{}此外克娄巴特拉与奥古斯都联合统治……十三年。\newline{}克娄巴特拉之后，奥古斯都又统治了……四十三年。\newline{}因为奥古斯都帝国的全部年数是……五十六年。\par

此外，让我们看看，在奥古斯都帝国第四十一年，当他（奥古斯都）在克利奥帕特拉死后统治了二十八年之后，基督是如何诞生的。（同一位奥古斯都在基督诞生后又活了十五年；而从克利奥帕特拉之死到基督诞生之日的剩余年数将把我们带到第四十一年，即奥古斯都在克利奥帕特拉死后第二十八年。）那么，到基督诞生之日，总共是三百三十七年零五个月：（由此充满了六十二个半七周：即四百三十七年零六个月。）然后，\Quote{永恒的正义}显现了，\Quote{至圣者被傅油}——即基督——并且\Quote{异象和先知\Emph{被封印}}，\Quote{罪过}被赦免，这些罪过通过信仰基督之名，为所有信祂的人所洗净。但祂说\Quote{异象和预言\Emph{被封印}}是什么意思？所有先知都曾预言祂要来并必须受苦。因此，既然预言通过祂的来临实现了，所以祂便说\Quote{异象和预言\Emph{被封印}}；因为祂是所有先知的印记，实现了他们昔日所预言的一切。因为在基督来临及其受难之后，不再有\Quote{异象或先知}预言祂的来临。简言之，如果不是这样，让犹太人在基督之后拿出任何先知的卷册、任何由天使所行的可见奇迹，（就像）过去列祖直到现今已来的基督来临之前所看到的那样；自那事件之后，\Quote{异象和预言被封印}，即被确认。福音作者恰当地写道：\Quote{法律和先知（止）于洗者若翰}。因为，在基督受洗时，即在祂自己的洗礼中圣化水时，过去所有属灵恩宠的满溢都在基督内停止了，祂封印了祂来临所实现的所有异象和预言。因此，祂最坚定地断言，祂的来临\Quote{封印了异象和预言}。\par

因此，显示了（正如我们所做的）年数，以及六十二个半七周完成的时间，在完成后，（我们已表明）基督已经来了，即已经诞生，让我们看看另外\Quote{七个半七周}，（这些七周是在前面七周被切割后细分出来的；让我们看看，即）它们是在什么事件中实现的：——\par

因为，在奥古斯都（他在基督诞生后存活了）之后，总共是……十五年（15）。\newline{} 继承他的是提庇留凯撒，统治了……二十年，七个月，二十八天（20等）。\newline{} （在他统治的第五十年，基督受难，当时祂约三十岁。）\newline{} 再次，加约凯撒，又称卡里古拉，……三年，八个月，十三天（3等）。\newline{} 尼禄凯撒，……十一年，九个月，十三天（11等）。\newline{} 加尔巴……七个月，六天（7等）。\newline{} 鄂图……三天。\newline{} 维泰留斯，……八个月，二十七天（8个月）。\par

维斯帕先在位第一年，在战争中征服了犹太人；总共是五十二年六个月。因为他统治了十一年。因此，在他们被攻陷的那一天，犹太人完成了达尼尔预言的七十个七周。\par

因此，当这些时间也完成，犹太人被征服后，此后\Quote{奠祭和祭献}在那地方停止了，从此以后不能再在那地方举行；因为\Quote{傅油}也在那地方\Quote{被消灭}了，在基督受难之后。因为曾预言傅油要在那地方被消灭；正如圣咏中所预言的：\Quote{他们刺穿了我的手和脚。}而这个\Quote{消灭}的苦难是在七十个七周的时间内，在提庇留凯撒、鲁贝留·杰米努斯和富菲乌斯·杰米努斯执政时期，在三月，逾越节期间，在四月朔日前第八天，无酵节的第一天，他们按照梅瑟所吩咐的，在晚上宰杀羔羊。因此，以色列全会众确实杀了祂，当比拉多想要释放祂时，对比拉多说：\Quote{祂的血归在我们和我们的儿女身上；}以及\Quote{你若释放祂，就不是凯撒的朋友；}\BibleRef{若 19:12}以便关于祂所写的一切得以应验。\par

% structure: p0036 source=h2 target=section reason=relative-heading-rank; base=h2

\section{第九章 论基督诞生与功绩的预言}

因此，我们开始证明基督的诞生是由先知预言的；如依撒意亚（例如）预言说：\Quote{达味家，请听！你们与人争论岂非小事？因为天主正在提出挑战。因此上主自己要给你们一个征兆：请看，贞女要怀孕生子，你要给祂起名叫厄玛奴耳}（解释为\InnerQuote{天主与我们同在}）：\Quote{祂要吃奶油和蜂蜜；}\Quote{因为，在孩子学会叫父亲或母亲之前，祂要得到大马士革的权势和撒玛利亚的战利品，与亚述王对抗。}\par

因此犹太人说：让我们挑战依撒意亚的预言，让我们进行比较，看是否在已经来临的基督身上，首先适用依撒意亚所预言的名称，以及（其次）祂宣布的关于祂的征兆。\par

那么，依撒意亚预言祂应当被称为厄玛奴耳；随后祂将夺取大马士革的权力和撒玛黎雅的掠物，对抗亚述王。他们说：\Quote{如今，你们的这位基督既没有被称为那名字，也没有从事战争。}但我们却认为，应当提醒他们回忆这段经文的 \Emph{上下文}。因为紧接着就是厄玛奴耳的解释——\Quote{天主与我们同在}——为使你们不仅注意名称的发音，也注意其意义。希伯来语的发音是厄玛奴耳，其解释是\Quote{天主与我们同在}。那么请问，自从基督的光照耀以来，\Quote{天主与我们同在}（即厄玛奴耳）这个说法是否普遍应用于基督？我想你们不会否认。因为那些从犹太教信仰基督的人，自从他们相信祂以来，每当他们想说厄玛奴耳时，都表示天主与我们同在；因此，一致认为那位曾被预言为厄玛奴耳的基督已经来临，因为厄玛奴耳所代表的已经来临——即\Quote{天主与我们同在}。同样，他们也被名称的发音所引导，将\Quote{大马士革的权力}、\Quote{撒玛黎雅的掠物}和\Quote{亚述王国}理解为好像预言基督是一位战士；他们没有注意到经文的前提：\Quote{因为孩子尚未知道叫父亲或母亲之前，必接收大马士革的权力和撒玛黎雅的掠物，对抗亚述王。}首先要看祂 \Emph{年龄} 的证明，看看那里所示的年龄是否可能表明基督已是一个 \Emph{成人}，更不用说 \Emph{将军} 了。的确，婴儿的哭声会召唤人们拿起武器，不是用号角而是用拨浪鼓发出战争信号，不是从马背或城墙上指出敌人，而是从哺乳者和乳母的背上或脖子上指出敌人，从而代替乳房征服大马士革和撒玛黎雅。（除非在你们中，婴儿冲入战场——我想，他们先涂油晾干，然后背着书包，以黄油为口粮，他们知道如何刺伤而不是如何撕裂胸脯！）当然，如果自然不允许这样——即在成年之前当兵，在认识父亲之前夺取\Quote{大马士革的权力}——那么这预言显然是有比喻性的。但他们又说：\Quote{但自然不允许\InnerQuote{童贞女}成为母亲；然而先知必须被相信。}这很合理，因为他为一件难以置信的事呼求信任，说这将是一个 \Emph{征兆}。\Quote{因此，}他说，\Quote{将给你们一个征兆。看，童贞女将怀孕生子。}但出自天主的征兆，除非包含某种奇特的新奇，否则就不会显得是征兆。总之，当你们急切地想要打倒任何人对这神圣预言的信仰，或转化那些单纯的人时，你们胆敢撒谎，仿佛经文包含\Quote{不是\InnerQuote{童贞女}，而是\InnerQuote{年轻女子}将怀孕生子}；你们甚至被这样的事实驳倒：日常发生的事——即年轻女子的怀孕和分娩——不可能显得是 \Emph{征兆}。因此，向我们呈现一位童贞母亲理应被相信是一个 \Emph{征兆}；但战士婴儿却不同。因为在这里不再涉及\Quote{征兆}的问题；而是，在授予了新奇诞生的\Quote{征兆}之后，紧接着的步骤是宣告婴儿将吃\Quote{蜂蜜和黄油}的不同后续安排。当然，这并非\Quote{征兆}；这是婴儿的天性。但祂将接收\Quote{大马士革的权力和撒玛黎雅的掠物，对抗亚述王}——这是一个奇妙的 \Emph{征兆}。坚持婴儿年龄的界限，探究预言的意义；不，不如将你们不愿相信的真相偿还，预言便通过其应验的关系变得可理解。请相信那些东方的贤士，他们以黄金和乳香供奉基督的婴儿时期，如同君王；婴儿没有战争和武器就接收了\Quote{大马士革的权力}。因为，除了众所周知东方的 \Emph{权力}——即 \Emph{力量}——通常以黄金和香料丰富之外，确定的是神圣圣经视 \Emph{黄金} 为构成其他民族 \Emph{权力} 的东西，正如它通过匝加利亚所说：\Quote{犹大将在耶路撒冷守护，并聚集周围所有民族的精力，黄金和白银。}因为这一 \Emph{黄金} 的礼物，达味也说：\Quote{祂将得到阿拉伯的黄金}；又说：\Quote{阿拉伯和舍巴的君王将给祂献上礼物。}因为东方通常认为贤士是君王；而大马士革在叙利亚分裂之前曾被认为属于阿拉伯，后来才转入叙利亚腓尼基：基督那时\Quote{接收}了该地的\Quote{权力}，即接收了其标志——黄金和香料。此外，祂也接收了\Quote{撒玛黎雅的掠物}，即贤士们本身，他们因引导和指示的星辰而认出祂，以礼物尊崇祂，并屈膝朝拜祂为主为王，从而成了\Quote{撒玛黎雅的掠物}，即偶像崇拜的掠物——因为他们相信了基督。因为（圣经）以\Quote{撒玛黎雅}之名指代偶像崇拜，撒玛黎雅因偶像崇拜而可耻；因为她当时在耶罗波安王领导下背叛了天主。这并非神圣圣经的新事，即基于 \Emph{罪行} 的相似而比喻性地使用 \Emph{名称} 转移。因为它称你们的首领为 \Emph{索多玛的首领}，你们的民众为 \Emph{哈摩辣的民众}，\BibleRef{依 1:10}，尽管那些城市早已消失。别处它通过先知对以色列民众说：\Emph{你们的父亲是阿摩黎人，你们的母亲是赫特人}；他们并非由那族所生，而是因不敬虔的相似而被称作他们的儿子，而天主曾通过依撒意亚先知称他们为 \Emph{自己的儿子}：\Quote{我养育并高举了儿子。}同样，埃及有时在那先知中因迷信和诅咒而被理解为整个世界。同样，巴比伦在我们的若望书中是罗马城的预象，因为同样伟大而骄傲于其统治，并战胜圣徒。因此，相应地，（圣经）也给贤士们冠以 \Emph{撒玛黎雅人} 的称号—— \Emph{被剥夺} 了他们曾与撒玛黎雅人共同拥有的东西，即我们所说的——偶像崇拜——对抗上主。（它说）\Emph{对抗}，此外，\Emph{亚述王}——对抗魔鬼，它至今仍以为自己统治着，如果它将圣徒从天主宗教中逐出的话。\par

此外，这种解释将得到支持，因为我们发现其他地方圣经也将基督描述为战士，正如我们从某些武器名称和这类词语中所看到的。但通过其余含义的比较，犹太人将被驳倒。\BibleQuote{你束上腰，将剑佩在大腿上。}但关于基督，你上面读到了什么？\BibleQuote{你在世人中美丽超群；恩宠倾注在你的唇上。}然而，如果他在称赞那人的美丽超群和唇上的恩宠，同时又让他佩剑备战，这是非常荒谬的；关于这人，他接下去说：\BibleQuote{愿你勇往直前，百战百胜，前进为王！}他又补充说：\BibleQuote{因为你的温和与正义。}有谁会使用剑而不实践温和与正义的反面，即诡诈、粗暴和不义，这些自然是战争事务的特点？那么，我们是否看到，那具有不同作用的是另一把剑——即\OriginalTerm{天主的圣言}{Divine word of God}，它被旧约和新约这两部约版双刃磨利，被它自身智慧的正义磨利，按各人的行为回报各人？因此，在圣咏中，天主基督可以不经战功而佩带天主圣言的象征性剑，这是合理的；所预言的\BibleQuote{美丽}和\BibleQuote{唇上的恩宠}与这把剑相符；当天主基督被宣布将要按天主圣父的旨意降临时，在\Person{达味}眼中，祂那时正是\BibleQuote{将剑佩在大腿上}的。他说：\BibleQuote{你右手的伟大将引导你}——即属灵恩宠的德能，基督的认知由此推导出来。他说：\BibleQuote{你的箭是锋利的}——天主到处飞驰的诫命（箭）威胁着揭露每颗心，给每个良心带来刺痛和穿透：\BibleQuote{万民将仆倒在你脚下}——当然，是朝拜。因此，基督在战争和武器方面如此强大；祂将不仅夺取\Place{撒玛黎雅}的战利品，也将夺取所有民族的。承认祂的\BibleQuote{战利品}是象征性的，因为你们已经知道祂的武器是寓言的。因此，到目前为止，已经来临的基督并不是战士，因为\Person{依撒意亚}没有这样预言祂。\par

\Quote{但如果那被认为将要来临的基督不被称为耶稣，为什么已经来临的被称为耶稣基督？}他们说。那么，每个名字将在天主的基督中相遇，在祂身上也发现了\Quote{耶稣}这个称呼。要知道你们错误的习惯性特征。在任命\Person{梅瑟}的继承人的过程中，\Person{农}的儿子\Person{何西阿}肯定从他原来的名字改变，开始被称为\Person{耶稣}。当然，你们说。我们首先断言这是未来的预像。因为，由于耶稣基督要引入第二子民（由我们这些先前被遗弃在世界中的外邦人所组成）进入\BibleQuote{流奶流蜜}的福地（即进入永生，没有比这更甜美的）；而这件事必须发生，不是通过\Person{梅瑟}（即不是通过法律的训诫），而是通过\Person{若苏厄}（即通过新法律的恩宠），在我们用\BibleQuote{火石刀}行割礼之后（即用基督的诫命，因为基督在很多方面和预像中被预言为磐石）；因此，那被准备作为这圣事图像的人，是在主的名号的预像下被任命的，甚至被命名为耶稣。因为那曾对\Person{梅瑟}说话的那一位是圣子本人；祂也总是\Emph{被看见}。\BibleRef{户 12:5}因为\Person{天主圣父}，没有人看见过而能活着。因此一致同意，天主圣子亲自对\Person{梅瑟}说话，并对百姓说：\BibleQuote{看，我派遣我的天使在你——即百姓——面前，在路上保护你，领你进入我为你预备的地方：你要留心听从他，不可悖逆他；因为他不会逃脱你的注意，因为我的名在他身上。}\BibleRef{出 23:20}因为\Person{若苏厄}将要带领百姓进入福地，而不是\Person{梅瑟}。现在祂称他为\Quote{天使}，是因为他要成就的伟大奇迹的宏大（\Person{农}的儿子\Person{若苏厄}成就了这些奇迹，你们自己也读到），以及他先知宣布（即神圣旨意）的职务；正如圣神以圣父的身份说话，通过先知称基督的前驱\Person{若翰}为将来的\Quote{天使}：\BibleQuote{看，我派遣我的天使在你——即基督——面前，他要在你面前预备你的道路。}圣神称那些祂委派作为祂权能仆人的为\Quote{天使}，这并非新做法。因为同样的\Person{若翰}不仅被称为基督的\Quote{天使}，而且被称为在基督面前照耀的\Quote{明灯}：因为\Person{达味}预言：\BibleQuote{我为我的基督预备了明灯；}\Person{基督}本人，来\Quote{满全先知}，也这样称呼\Person{若翰}对犹太人。\BibleQuote{祂是燃亮的明灯，}祂说，因为他不仅\Quote{在旷野预备祂的道路}，而且通过指出\BibleQuote{天主的羔羊}，照亮了人们的心灵，使他们明白祂就是\Person{梅瑟}常宣布要受苦的那只羔羊。同样，（\Person{农}的儿子被称为）\Person{若苏厄}，是因为他名字未来的奥秘：因为那对\Person{梅瑟}说话的那位（祂）确认了祂自己所赐给他的名字为自己的名字，因为祂命令他从那时起不再被称为\Quote{天使}或\Quote{何西阿}，而是\Quote{若苏厄}。因此，每个名字都适合天主的基督——祂也应被称为耶稣（正如基督）。\par

至于那当从她而生\Person{基督}（如上文所述）的童贞女，必须追溯她的谱系于\Person{达味}的后裔，先知在随后的段落中明确断言：\BibleQuote{将有一枝条从\Person{叶瑟}的根生出}——这枝条就是\Person{玛利亚}——\BibleQuote{并从他的根上发出花朵：上主的神，智慧和聪敏的神，明达和虔诚的神，超见和真理的神，将要降在祂身上；上主敬畏的神必充满祂。} 因为普世性的属灵凭证之总和，除\Person{基督}外，无人当得起；祂因荣耀、因恩宠而被比作\Quote{花朵}；但被称为\Quote{出于叶瑟的根}，祂的起源由此追溯——即经由\Person{玛利亚}。因为祂生于\Place{白冷}的本土，出于\Person{达味}的家；如同在\Place{罗马人}中间，\Person{玛利亚}在户口登记中被记载，而\Person{基督}由她所生。\par

我再次要求——承认那屡次被预言将出于\Person{叶瑟}族裔的那位，将显出全部谦卑、忍耐和宁静——祂来了吗？同样（在此如在前者），那显出此种品格的人将是已经来临的\Person{基督}本身。因为关于祂，先知说：\BibleQuote{他备受痛苦，熟悉病弱}；祂\BibleQuote{被牵去如同待宰的羊，又如羔羊在剪毛者前缄口不语。} 如果祂\BibleQuote{既不争辩，也不呼喊，街上听不到他的声音}；祂\BibleQuote{不折断压伤的芦苇}——即以色列的信德——\BibleQuote{也不熄灭将残的灯火}——即外邦人短暂的光辉——反而因自己真光的升起使之更加明亮——那么祂只能是那被预言的一位。因此，已来临的\Person{基督}的行为，必须对照圣经的准则来审查。因为，若我没弄错，我们看见祂有双重作为——即\Emph{宣讲}与\Emph{德能}。现在，让我们简要处理每一项。因此，让我们按已定下的顺序展开，教导\Person{基督}被宣告为一位\Emph{宣讲者}；如通过\Person{依撒意亚}所说：\BibleQuote{你大声呼喊，不要停止，提高你的声音，如同号角，向我的百姓宣告他们的过犯，向雅各伯家宣告他们的罪恶。他们天天寻求我，乐意学习我的道路，好像行义未曾离弃天主判断的百姓，}等等。此外，祂要从\Person{圣父}那里施行\Emph{德能}：\BibleQuote{看哪，我们的天主要施行报应；祂要亲自来拯救我们。那时，瞎子要看见，聋子要听见，哑巴的舌头要歌唱，瘸子要跳跃如鹿，}等等。这些工作连你们也不否认\Person{基督}行了，因为你们常说：\Quote{为了这些工作，你们没有用石头砸祂，是因为祂在安息日做了这些事。}\par

% structure: p0044 source=h2 target=section reason=relative-heading-rank; base=h2

\section{第十章 论\Person{基督}的苦难及其在旧约中的预言与预象}

关于祂苦难的最后一步，显然，你们提出质疑；断言十字架的苦难并未预言关乎基督，并极力主张，天主将祂自己的独生子暴露于那种死亡方式是不可信的；因为祂自己曾说：\BibleQuote{凡挂在树上的，都是被咒骂的。} 但事情的原因先已解释了这咒骂的含义；因为祂在申命纪中说：\BibleQuote{若有人犯了该处死的罪，被处死，你将他挂在树上，他的尸体不可留在树上过夜，务必当日将他埋葬，因为被挂在树上的人是天主所咒骂的；你不可玷污上主你的天主所赐给你作产业的地。} 因此，祂并非以咒骂之意判定\Person{基督}受此苦难，而是加以区分：凡\Emph{因任何罪}而犯了该处死的判决，且死在树上被挂起的人，他才是\BibleQuote{被天主咒骂的}，因为是他自己的罪导致他被挂在树上。相反，\Person{基督}口中没有诡诈，行事全然公义和谦卑，不仅（如上文所记所预言的）不是\Emph{因自己的功过}而受那种死亡，而是为使先知所预言将借你们之手临到祂身上的话得以应验；正如在圣咏中，\Person{基督}的圣神早已歌唱说：\BibleQuote{他们以恶报善；} \BibleQuote{我未曾抢夺的，如今却要偿还；} \BibleQuote{他们钉穿了我的手和脚；} \BibleQuote{他们给我的食物掺了苦胆，我渴时，他们拿醋给我喝；} \BibleQuote{他们为我的衣服拈阄；} 正如你们要在祂身上犯的其他（暴行）也都预被言明——凡此种种，祂真实而全然忍受，不是因祂自己任何恶行而受苦，而是\BibleQuote{为使先知口中的圣经得以应验。}\par

当然，苦难本身的奥秘本应在预言中以预象方式预先表明；这奥秘（愈不可思议），若被赤裸裸地预言，就愈可能成为\BibleQuote{绊脚石}；愈是宏伟，就愈需要被\Emph{预象}，以使理解其困难者可寻求天主恩宠的（帮助）。\par

因此，首先，依撒格被父亲带去作为祭品，他自己背负着\BibleQuote{木柴}，早在那个时代就指向了\Person{基督}的死亡；祂由圣父交出作为祭品，正如祂背着祂自己苦难的\BibleQuote{木柴}。\par

再者，若瑟本人也唯独在这一点上成为基督的预象（不再多提，以免延误我的进程），即他因天主的宠爱而遭受弟兄们的迫害，并被卖到埃及；正如基督被以色列出卖——（因此，）\Quote{按肉体说}，被祂的\Quote{弟兄}\BibleRef{罗 9:5}——当祂被犹达斯背叛的时候。因为若瑟也蒙父亲以这样的形式祝福：\Quote{他的荣耀（是）公牛的荣耀；他的角，是独角兽的角；他要借此将万民抛掷，直到地极。}当然，那里所指的并非一只独角的犀牛，也非双角的弥诺陶洛斯。而是基督被预示了：\Quote{公牛}，因祂两种性格中的每一种——对有些人严厉，如同审判者；对另一些人温柔，如同救主；祂的\Quote{角}将是\Emph{十字架}的两端。因为即使在船桁——十字架的一部分——中，两端也以此名称呼；而桅杆的中柱则是\Quote{独角兽}。事实上，藉着十字架的这一权能，并如此带着角，祂现在一方面藉着\Emph{信德}\Quote{抛掷}万民，将他们从大地吹向天堂；另一方面，将来有一天，祂要藉着\Emph{审判}\Quote{抛掷}他们，将他们从天堂掷向大地。\par

祂，再次，在另一处同一经文中将是那\Quote{公牛}。当雅各伯为西默盎和肋未祝福时，他预言了经师和法利塞人；因为他们的起源出于此。因为（他的祝福）按属灵的意义如此解释：\Quote{西默盎和肋未因他们的派别成就了不义}——由此，他们迫害了基督：\Quote{我的灵魂不要进入他们的计谋！我的心不要停留在他们的会集！因为他们在愤怒中杀了人}——即先知——\Quote{并在他们的贪欲中砍断了公牛的腿筋！}——即基督，祂——在先知被杀之后——被他们杀害，并用钉子刺穿祂的筋腱，耗尽了他们的残暴。否则，如果他们已犯了谋杀罪，他却指责别人而不指责他们屠戮，那便是无意义的。\par

但是，现在来看梅瑟，我想知道，为什么他只是在若苏厄与阿玛肋克交战时，才\Emph{坐着}伸开双手祈祷，而在如此危急的情况下，他本当更应该屈膝、捶胸、面伏于地来呈上他的祈祷；除非是因为那里，主耶稣的名号是谈论的主题——祂注定将来要单独与魔鬼交锋——而\Emph{十字架}的形像也是必要的，（那个形像）耶稣要藉此赢得胜利？同样，为什么同一位梅瑟，在禁止任何\Quote{偶像}之后\BibleRef{出 20:4}，却铸造了一条铜蛇，挂在\Quote{木杆}上，悬吊着，作为以色列得医治的景象，当时在他们拜偶像之后，正被蛇咬而遭灭绝，除非在此他是在展示主的\Emph{十字架}，其上那\Quote{蛇}——魔鬼——被\Quote{展示出来}，并且为每一个被这种蛇（即他的天使）咬伤的人，从罪愆中彻底回转，归向基督\Emph{十字架}的圣事，救恩便得以成就？因为那时注视那（十字架）的人，便从蛇咬中被解救。\par

现在，如果你在圣咏中先知的话里读到：\Quote{天主从\Emph{树木}之上为王}，我等着听你对此如何理解；恐怕你或许以为这里指某个木匠国王，而不是基督——祂从那一刻起为王，当祂克服了因\Quote{树木}的苦难而来的死亡。\par

同样，依撒意亚又说：\Quote{因为有一个婴孩为我们诞生了，有一个儿子赐给了我们。}这有什么新奇之处，除非他指的是天主的\Quote{儿子}？——且有一个婴孩为我们诞生，祂统权的起始已放在\Quote{祂的肩上}。世上哪个国王将权力的标志戴在\Emph{肩上}，而不在头上戴冠冕，或手中持权杖，或有什么特别的服饰标记？但那位\Quote{万世之君}、\Emph{新颖}的基督耶稣，唯独将祂\Emph{崭新}的荣耀、权能和崇高——即\Emph{十字架}——\Quote{放在祂的肩上}；这样，按照先前的预言，主从此\Quote{从树木之上为王}。因为关于这棵树，天主也藉耶肋米亚暗示，你会说：\Quote{来，让我们把\Emph{木头}放在他的\Emph{饼里}，并把他从活人之地抹去；他的名字不再被记念。}当然，那\Quote{木头}是放在祂的\Emph{身体}上；因为基督如此启示了，称祂的身体为\Quote{饼}，先知在古时以\Quote{饼}一词宣告了祂的身体。如果你还要寻求关于主十字架的预言，第二十一篇圣咏最终必能使你满足，因为它包含了基督的全部苦难；祂甚至在那样早的时候，就歌唱自己的荣耀。\Quote{他们刺穿了}，祂说，\Quote{我的手和我的脚}——这是十字架特有的残酷；而当祂恳求天父的助佑时，又说：\Quote{拯救我}，\Quote{脱离狮子的口}——当然是指死亡——\Quote{并使我卑微的身体脱离独角兽的角}——即十字架的两端，如上文所示；这十字架，大卫本人未曾受，犹太的众王也未曾受：这样你就不会认为这里所预言的是其他某个人的苦难，而是那唯独被百姓如此显著地钉在十字架上者的苦难。\par

现在，如果你心中的硬顽坚持拒绝和嘲笑所有这些解释，我们将证明：只要基督的\Emph{死亡}曾被预言就够了，因为既然死亡的性质未曾被指定，就可以理解它是藉着\Emph{十字架}而实现的，并且\Emph{十字架}的苦难不应归给任何人，只应归给那位其\Emph{死亡}不断被预言者。因为我想在依撒意亚的一句话中显示祂的\Emph{死亡}、\Emph{苦难}和\Emph{埋葬}。他说：\Quote{因我百姓的罪恶，祂被引至死亡；我要以恶人作为祂的埋葬，以富人作为祂的死亡，因为祂没有行过不义，口中也没有诡诈；天主愿意救赎祂的灵魂脱离死亡，}等等。他又说：\Quote{祂的埋葬已从中被取走。}因为祂除非死了，否则不会被埋葬；除非藉着祂的复活，否则祂的埋葬不会从中被取走。最后他补充说：\Quote{因此祂必得许多人为产业，也必分取许多人的掠物：}除了那位\Quote{被生下来}的（如我们以上所已指出的），还有谁呢？——\Quote{因为祂的灵魂被交付于死亡？}因为，既然祂蒙受恩宠的\Emph{原因}已被显示——即为了补偿必须补偿的\Emph{死亡}之伤害——同样也显示，祂注定要因死亡而获得这些赏报，是要在死亡\Emph{之后}获得——当然是在\Emph{复活}之后。因为祂受难时所发生的事，即正午变为黑暗，亚毛斯先知宣告说：\Quote{他说：\InnerQuote{在那一天，主说：太阳要在正午落下，光明之日要在以上变为黑暗；我要将你们的庆节变为悲哀，将你们的歌曲变为哀歌；我要把麻布束在你们腰间，把剃头放在你们头上；我要使悲哀如同丧独生子，使与他同在的人如同哀悼之日。}}因为你们要在你们新年的第一月初这样做，连梅瑟也曾预言过他预告以色列全体子民要在黄昏时宰杀一只羔羊，并要带着苦菜吃这日的隆重祭献（即无酵饼的逾越节）；并加上说：\Quote{这就是\Emph{主的逾越节}，}即\Emph{基督的苦难}。这预言如此应验了：你们在\Quote{无酵饼的第一日}宰杀了基督；\BibleRef{格前 5:7}并且（为使预言应验）那日急忙成为\Quote{黄昏}——即造成黑暗，这黑暗在正午发生；因此\Quote{天主将你们的庆节变为悲哀，将你们的歌曲变为哀歌。}因为在基督受难之后，你们甚至被掳掠和分散，这早已藉着圣神预言过了。\par

% structure: p0054 source=h2 target=section reason=relative-heading-rank; base=h2

\section{第十一章 来自厄则克耳的进一步证明 先知论证迄今的总结}

再者，正是因你们的这些罪行，厄则克耳宣布你们的毁灭即将来临：不仅在此世代——这毁灭已降临——亦在随后的\BibleQuote{报应的日子，}\BibleRef{依 61:2}中。从这毁灭中，唯有一人得获释放，即那额上被你们所弃的基督的苦难所印封的人。因为经上这样记载：\BibleQuote{主对我说：\InnerQuote{人子，你看见以色列的长老们在做什么？各人在黑暗中，各人在隐密的卧室里，因为他们说：\InnerQuote{主看不见我们，主已经遗弃了大地。}}主又对我说：\InnerQuote{你再去，要看见他们做更大的恶事。}他领我到上主殿宇朝北的门限旁；看，那里有妇女坐着哀悼塔模斯。主对我说：\InnerQuote{人子，你看见了吗？犹大之家做这些恶事还算\Emph{轻微}吗？但你还要看见他们更可憎的事。}他领我进入上主殿宇的内院；看，在上主的殿门口，在门廊和祭台之间，约有二十五人，背向上主的殿，面向东方；他们正在朝拜太阳。主对我说：\InnerQuote{人子，你看见了吗？犹大之家做这些事算小事吗？因为他们恶贯满盈，看，\Emph{正是}他们自己仿佛作鬼脸；我必发怒行事，我的眼不顾惜，也不可怜；他们必大声呼号在我的耳中，我却不听他们，也不可怜。}他大声向我耳中呼喊说：\InnerQuote{这城的惩罚近了；各人手中都拿着毁灭的器械。}看，有六人从朝北的上门那边走来，各人手中拿着散播的斧头；他们中间有一个人身穿长衣，\BibleRef{默 1:13}腰束蓝宝石带；他们进入，站在铜祭台旁。以色列天主的荣耀曾停在殿上，在殿的院内，从革鲁宾上升起；主呼唤那身穿长衣、腰间束带的人说：\InnerQuote{你走遍耶路撒冷，在那些为城中一切可憎之事而叹息哀哭的人额上划十字记号。}他正做着这些事时，对那听者说：\InnerQuote{你跟着他进城，要砍断，你的眼不要顾惜，也不要可怜老人、青年或童女；要杀死所有的小孩子和妇女，使他们彻底消灭；但额上有十字记号的人不可接近；先从我的圣者开始。}}这\Quote{记号}的奥秘曾以多种方式预言；这\Quote{记号}为人类预置了生命的基础；这\Quote{记号}是犹太人不会相信的：正如梅瑟从前在\WorkTitle{出谷纪}中不断宣告说：\BibleQuote{你会被驱逐出你要进入的那地；你在那些民族中不得安息；你的脚掌不得安稳；天主必使你心疲力倦，灵魂消沉，眼睛昏花，不能看见；你的生命必悬在树上，悬在你眼前；你也不得信赖你的生命。}\par

因此，既然预言藉祂的来临——即藉我们上面所纪念的诞生和我们已清楚解释的苦难——而实现了，这也就是达尼尔说\Quote{异象和预言\Emph{被封闭了}；}的原因；因为基督是众先知的\Quote{印玺}，成全了古时关于祂所宣告的一切：因为自从祂来临和亲自受苦之后，便不再有\Quote{异象}或\Quote{先知}；因此他极强调地说，祂的来临\Quote{\Emph{封闭}了异象和预言。}这样，藉显示\Quote{年数的数目，以及满了六十二个半星期的时期}，我们证明在指定的时间基督来了，即降生了；又（藉显示）那\Quote{七个半星期}（这些星期被分割，从先前的星期中截断出来）的期间，我们证明基督受了苦难，并藉着随之而来的\Quote{七十个星期}的结束和城市的毁灭，（我们证明）\Quote{祭献和傅油}从此止息。\par

至此，在这些事上，我们已经足够跟踪了基督所命定的道路，由此证明祂正如从前所宣告的，即使是基于圣经的一致性，这使我们能够根据多数人的预判，起来反对犹太人。因为他们不可质疑或否认我们所引用的著作；事实上，那些预言在\Emph{基督之后}注定要发生的事，如今已被承认为应验了，这使他们不可能否认这些著作与神圣经书相等。否则，除非祂已经来临——\Emph{在祂之后}那些常被宣告的事必须成就——否则那些已经完成的事怎能得证呢？\par

% structure: p0058 source=h2 target=section reason=relative-heading-rank; base=h2

\section{第十二章　由外邦人的蒙召而得的进一步证据}

请看看那些自此从人类错误的旋涡中涌现出来的普世万民，归向主天主创造主和他的基督；如果你敢否认这是预言过的，那么圣父在圣咏中的应许立刻浮现在你眼前，其中说：\Quote{你是我的儿子，我今日生了你。你向我请求，我必将外邦人赐你作基业，将地极赐你作田产。}因为你无法肯定那\Quote{儿子}是指达味而非基督；也无法肯定那\Quote{地极}是应许给只在犹大单一地境作王的达味，而非应许给那已以其福音的信德掳获全世界的基督；正如他通过依撒意亚所说：\Quote{看，我已立你作我家族的盟约，作外邦人的光，为开启盲人的眼睛}——当然是指那些迷途的人——\Quote{释放被捆锁的}——即从罪恶中释放——\Quote{并从监狱之屋}——即死亡——\Quote{拯救坐在黑暗中的人}——即无知。如果这些福泽经由基督而来，那么它们就不是预言另一位，而是正预言那位我们视之为成就这一切的那位。\BibleRef{路 2:25}\par

% structure: p0060 source=h2 target=section reason=relative-heading-rank; base=h2

\section{第十三章　从耶路撒冷的毁灭和犹大地的荒凉提出的论证}

因此，既然以色列的子孙断言我们接受那已来临的基督是错误，就让我们根据圣经本身向他们提出抗辩：那被预言的基督确实已经来临；尽管借着达尼尔预言的时间，我们已经证明那被宣告的基督已经来临。现在，他理应在犹大的白冷出生。因为先知这样写着：\Quote{而你，白冷，在犹大的首领中决不是最小的，因为将由你出来一位领袖，他将牧养我的百姓以色列。}但如果直到如今他仍未出生，那么那被预告要从犹大支派中、从白冷出来的\Quote{领袖}又在哪里？因为他理应从犹大支派、从白冷出来。然而我们看见，如今没有一个以色列的子孙留在白冷；自从颁布禁令禁止任何犹太人在该地区边界停留以来，就一直如此，为要完全应验这一预言：\Quote{你的土地荒凉，你的城邑被火焚烧，}——即他预言在战争时期他们将遭遇的事——\Quote{你的地域在你眼前被外邦人吞吃，并且荒凉，被外族人民倾覆。}在另一处，借着先知如此说：\Quote{你要看见君王带着他的荣耀}——即基督，在天主圣父的光荣中行奇事；\Quote{你的眼睛要远远看见那地，}\BibleRef{依 33:17}——这正是你们所做的，由于你们的罪过，在耶路撒冷被攻陷后，你们被禁止进入你们的土地；只准许你们从远处用眼睛看见它：\Quote{你的心，}他说，\Quote{要默想恐怖，}\BibleRef{依 33:18}——即在他们遭受自身毁灭的时候。那么，一个\Quote{领袖}怎能从犹大出生，又怎么能从白冷\Quote{出来}，正如先知的神圣经卷所明确宣布的那样？因为今天那里已经没有留下一个以色列（家族）的人，基督能从其血统诞生呢？\par

现在，如果（按照犹太人的说法）他至今尚未来临，那么当他开始来临时，他将从何处受傅？因为法律规定，在掳掠中，不得调制王权的圣油之傅。\BibleRef{出 30:22}但若那里不再有\Quote{傅油}，正如达尼尔所预言的（因为他断言：\Quote{傅油将被铲除}），那么他们就不再拥有它，因为他们也没有了\Quote{圣殿}，那里曾有君王受傅所用的\Quote{角}。既然如此，就没有傅油，那将在白冷出生的\Quote{领袖}将从何处受傅？或者他怎能从白冷\Quote{出来}，既然在白冷根本没有以色列的后裔存在呢？\par

事实上，让我们第二次证明基督已经来临（正如先知所预言的），已经受难，已经回到天上，并要照预言所预定的从那里再来。因为按\Person{达尼尔}所言，我们读到，在祂来临之后，那城本身必须被毁灭；我们认出这事已经发生。因为圣经如此说：\Quote{城与圣所必与\Emph{领袖}一同被毁灭}\BibleRef{达 9:26}，——无疑就是那要从\Quote{白冷}、从\Quote{犹大}支派来的\OriginalTerm{领袖}{leader}。由此可见，当它的\Quote{领袖}必须在城中受难时，\Quote{城必同时被毁灭}，正如先知们的圣经所预言，他们曾说：\Quote{我终日向那悖逆抗拒我的百姓伸出双手，他们行走不善良的道路，只随从自己的罪恶。}在圣咏中，\Person{达味}说：\Quote{他们刺穿了我的手和脚，数算了我的所有骨头；他们又注视观看我，我渴时，他们拿醋给我喝。}这些事\Person{达味}并未遭受，以致他看似合理地说到自己；而是那被钉十字架的基督。况且，\Quote{手和脚}除了被悬在\Quote{\OriginalTerm{木}{tree}}上的那一位，是不会\Quote{被毁}的。因此，\Person{达味}又说：\Quote{主将从\Emph{木上}为王}：因为其他地方，先知也预言这\Quote{\OriginalTerm{木}{tree}}的果实，说\Quote{大地已给出了她的祝福}，——当然是指那童贞之地，尚未被雨水灌溉，也未被阵雨滋润，男人昔日从那里被首先塑造，而现在基督藉着肉身从童贞女从那里出生；\Quote{而\Emph{那木}，}他说，\Quote{结出了它的果实，}——不是乐园中那给原祖带来死亡的\Quote{树}，而是基督苦难的\Quote{\OriginalTerm{木}{tree}}，生命悬于其上，却被你们不信！因为这\Quote{木}在奥秘中，昔日\Person{梅瑟}曾用以使苦水变甜；由此，在旷野中因口渴而灭亡的百姓，喝了便复活了；正如我们一样，我们从外邦的灾难中被救出，我们曾停留其中因口渴而灭亡（即被剥夺天主圣言），喝了藉着\Quote{信德，这信德是在祂内}的基督苦难之\Quote{木}的洗礼水，便复活了——以色列却失落了这信德，如\Person{耶肋米亚}所说：\Quote{你们去问一问，是否发生过这样的事，民族能否更换自己的神祇（而它们并不是神！）。但我的子民更换了自己的光荣：因此对他们毫无益处：天为此惊恐}（何时惊恐？无疑是在基督受难时），\Quote{且极其战栗，}他说，\Quote{且正午太阳昏暗}（何时\Quote{极其战栗}，除非在基督受难时，那时地也震动到中心，圣所帐幔撕裂，坟墓崩开？\Quote{因为我的子民做了两件恶事：我，}祂说，\Quote{他们完全离弃了，活水的泉源，又为自己凿了破裂的水池，不能蓄水。}）无疑，由于不接收基督，\Quote{活水的泉源}，他们开始拥有\Quote{破裂的水池}，即供\Quote{外邦人分散}使用的会堂，圣神不再停留其中，如同往日基督来临前祂惯常在圣殿中停留一样，基督才是天主的真圣殿。因为，他们还要遭受这种对天主圣神的渴望，先知\Person{依撒意亚}曾说过：\Quote{看哪，事奉我的要吃喝，你们却要饥饿；事奉我的要喝，你们却要口渴，并因精神上的普遍苦难而哀号：因为你们要把你们的名字留作我的选民饱足之用，但上主必将你们杀死；而事奉我的将得一个新名，这名字在列邦中必受祝福。}\par

再者，这\Quote{\OriginalTerm{木}{tree}}的奥秘，我们在\WorkTitle{列王纪}中也读到被颂扬。因为当先知弟子们在约但河岸用斧头砍\Quote{\OriginalTerm{木}{wood}}时，铁头脱落并沉入水中；于是，当先知\Person{厄里叟}上来时，先知弟子们求他从水中取出沉下的铁。\Person{厄里叟}于是取了\Quote{\OriginalTerm{木}{wood}}，投入铁器沉没之处，铁立刻浮起漂在水面，而\Quote{木}却沉下，先知弟子们将其拾回。由此他们明白\Person{厄里亚}的神灵已赐予他。还有什么比这\Quote{木}的奥秘更显明呢？——这世界的顽硬曾沉于错误的深渊，藉着基督的\Quote{\OriginalTerm{木}{wood}}，即祂的苦难，在洗礼中得释放；这样，先前在亚当内因\Quote{\OriginalTerm{树}{tree}}而丧亡的，应在基督内藉\Quote{\OriginalTerm{木}{tree}}得以恢复。而我们当然，继承并占据了先知的位置，今日在世上承受了先知们因神圣宗教所受的对待：因为他们有的被石头砸死，有的被放逐；但更多的，则被交付于死亡杀戮——这是他们无法否认的事实。\par

这\Quote{木料，}同样地，\Person{亚巴郎}之子\Person{依撒格}亲自背在自己的祭献上，当时\Person{天主}命令他把自己作为牺牲献给祂。然而，由于这些奥秘一直在保留，以便在\Person{基督}的时代完美应验，一方面，\Person{依撒格}和他的\Quote{木料}被保留了，而那只被灌木丛中的角卡住的公羊被献上；另一方面，\Person{基督}在祂的时代，将祂的\Quote{木料}背在祂自己的肩上，紧贴十字架的角，头上戴着荆棘冠冕。因为祂必须为所有外邦人成为祭献，祂\BibleQuote{被牵到宰杀之地，又像羔羊在剪毛者前无声，所以祂不开口}（因为祂在\Person{比拉多}审问祂时，一言不发）；因为\BibleQuote{祂的审判被夺去了，祂的出生，又有谁能述说呢？}因为根本没有人知道\Person{基督}在成胎时的诞生，当时\Person{童贞玛利亚}因天主的话而怀孕；并且因为\BibleQuote{祂的生命将从地上被夺去。}那么，为什么在祂从死者中复活（这在第三日实现）之后，诸天又接回祂呢？这是符合\Person{欧瑟亚}先知所说的：\BibleQuote{他们要在黎明前起来到我这里，说：\InnerQuote{让我们去，回到上主我们的天主那里，因为祂将牵引我们出来，释放我们。}过了两天，在第三天}——这是祂光荣的复活——诸天接回了祂（\Person{圣神}自己也曾从那里来到童贞女那里），祂的诞生和苦难，犹太人都没有承认。因此，既然犹太人仍然争辩说\Person{基督}尚未来临，而我们已在许多方面证实祂已来临，就让犹太人认清自己的命运——这是他们在\Person{基督}来临后因着鄙视和杀害祂的不敬而不断被预言将要承受的命运。因为首先，从\Person{依撒意亚}所说\BibleQuote{人抛弃了金银制成的可憎之物，他们制造它们用虚空有害的礼仪来朝拜}的那一天起——即自从我们外邦人，通过\Person{基督}的真理获得双重光照后，抛弃了我们的偶像（让犹太人看见）——以下的事也照样应验了。因为\BibleQuote{\OriginalTerm{万军的上主}{Lord of Sabaoth} \Emph{夺去了}，在耶路撒冷犹太人中间}其他提到的诸事中，也夺去了\BibleQuote{智慧的建筑师}，他建造教会、\Person{天主}的殿、圣城和上主的殿。因为从那时起，\Person{天主}的恩宠在她们中间停止了工作。并且\BibleQuote{云被命令不得降雨在\Place{索肋克}葡萄园上}——云是属天的恩惠，它们被命令不得降给以色列家；因为它\BibleQuote{长出了\Emph{荆棘}}——以色列家用此荆棘为\Person{基督}编了一顶冠冕——而不是\BibleQuote{\Emph{正义}，而是\Emph{喧嚷}}——这喧嚷就是他们强逼将祂交给十字架的声音。因此，先前的恩宠礼物被撤回了，\BibleQuote{法律和先知到\Person{若翰}为止}，以及\Place{贝特赛达}的鱼池直到\Person{基督}来临：此后它就不再医治性地除去以色列的健康疾病；因为他们坚持疯狂，\Person{主}的名通过他们被亵渎，正如经上所记：\BibleQuote{因你们的关系，\Person{天主}的名在外邦人中受亵渎}：因为那（附于那名）的恶名从他们开始，并在\Person{提庇留}到\Person{韦斯巴芗}的期间传播。因为他们犯了这些罪行，且不明白\Person{基督}\BibleQuote{将被找到}在\BibleQuote{他们被眷顾的时候}，他们的土地变成了\BibleQuote{荒芜，城市被火焚烧，陌生人当着他们的面吞吃他们的区域：\Place{熙雍}的女儿被遗弃，像葡萄园中的瞭望台，或黄瓜园中的茅棚}——自从\BibleQuote{\Place{以色列}不认识}\Person{主}，\BibleQuote{百姓不明白祂}；反而\BibleQuote{完全离弃，并激起以色列的圣者的愤怒}的时候。同样地，我们发现了一个\Emph{有条件}的\Emph{剑}的威胁：\BibleQuote{如果你们不愿意，不听从，剑将吞噬你们。}\BibleRef{依 1:20}由此我们证明\Emph{剑}就是\Person{基督}，因不听从祂而灭亡；祂在圣咏中又向圣父要求分散他们，说：\BibleQuote{用祢的能力分散他们}；此外，又通过\Person{依撒意亚}为他们的\Emph{彻底焚烧}祈祷。\BibleQuote{因我之故，}祂说，\BibleQuote{这些事临到你们；你们要在焦虑中安眠。}\par

因此，既然犹太人被预言注定要因\Person{基督}而遭受这些灾难，并且我们发现他们已经遭受了，并看到他们被分散并一直留于分散中，显然这些事是因\Person{基督}而降临到犹太人身上的，圣经的意义与事件的结果和时间的顺序相一致。否则，如果\Person{基督}尚未来临，因祂的缘故他们被预言将要遭受这些，那么当祂来临后，他们必然也会这样遭受。那时哪里还有\Place{熙雍}的女儿被遗弃，而她现在已不存在？哪里还有城市被烧毁，而它们早已成为灰烬和废墟？哪里还有种族的分散，而他们现在已在流徙？恢复\Person{基督}将要找到的\Place{犹大}地的状况，然后（如果你愿意）再争辩说有另一位（\Person{基督}）要来。\par

% structure: p0067 source=h2 target=section reason=relative-heading-rank; base=h2

\section{第十四章。结论。犹太人错误的线索。}

现在（除了直接的问题之外）学习你们错误的线索。我们确认，先知们所预言的基督有\Emph{两种}特性，并预告了他的\Emph{两次来临}：一次是在谦卑中（当然是第一次），那时他要被领去\BibleQuote{如同羊羔被牵去宰杀，又像母羊在剪毛者前缄默，同样他也没有开口}，甚至连容貌也不俊美。因为\BibleQuote{我们关于他所宣布的}，\Person{先知}说，\BibleQuote{他好像一个孩童，像干旱地上的一根苗；在他身上没有可爱或荣耀。我们看见了他，他没有可爱或恩宠；但他的面容是不光彩的，比不上世人。}\BibleQuote{他是一个\Emph{被置于}瘟疫中的人，知道怎样忍受软弱：}即被圣父立为\BibleQuote{绊脚石}，并被他\BibleQuote{稍微降于天使}，他自称\BibleQuote{是一条虫，不是一个人，是人的羞辱，和百姓的渣滓。}这些卑微的证据符合\OriginalTerm{第一次来临}{First Advent}，正如那些崇高的证据符合\OriginalTerm{第二次来临}{Second Advent}；那时他将不再成为\BibleQuote{绊脚石或绊倒的盘石}，而是\BibleQuote{最高角石}，在地被弃绝后（在地上）被举升（到天上）并高升以臻圆满，如\BibleRef{弗 1:10}，以及那块\BibleQuote{盘石}——我们必须承认——在\Person{达内尔}书中读到的从山上非人手凿出的石头，要粉碎并磨灭世俗王国的像。关于同一个（基督）的第二次来临，\Person{达内尔}曾说：\BibleQuote{看，仿佛一位人子，乘着天上的云彩，来到万古常存者面前，被引领到他面前。站在旁边的人领他到他跟前。赐给他王权；世上万民按各自的种族，以及一切荣耀，都要侍奉他：他的权柄\Emph{是}永远的，不能被夺去，他的国度\Emph{一个}不能朽坏的。}那时，他当然要有荣耀的面容和恩宠，不\BibleQuote{比世人更为欠缺}；因为（他将是）\BibleQuote{比世人更俊美}。\Person{圣咏作者}说：\BibleQuote{恩宠倾注在你口中：因此天主永远赐福了你。佩上你的剑于腰间，在荣耀和美丽中得胜！}而圣父后来使他稍微低于天使之后，\BibleQuote{以光荣和尊贵加冕于他，将\Emph{万有}置于他的脚下。}那时他们要\BibleQuote{认识他们所刺透的那一位，并要各支派捶胸痛哭}；当然是因为在往日，当他们处于人类卑微境况时，他们\Emph{不}认识他。\Person{耶肋米亚}说：\BibleQuote{他是个人，谁能认识他？}因为，\Person{依撒意亚}说：\BibleQuote{他的诞生谁能述说？}同样，在\Person{匝加利亚}中，在他自己的身份中，甚至在他名字的奥义中，圣父最真实的司铎，他自己的基督，以两件衣服描绘了\OriginalTerm{两次来临}{two advents}。首先，他穿着\BibleQuote{污秽的衣服}，即受难和必死肉身的卑微，当时魔鬼也反对他——即叛徒\Person{犹达斯}的唆使者——甚至在他受洗后曾试探他。其次，他被剥去先前污秽的衣服，穿上长袍，并戴上头巾和干净的冠冕，即（象征）\OriginalTerm{第二次来临}{Second Advent}；因为他被显示已获得了\BibleQuote{光荣和尊贵}。你不能说（那里描述的）那个人是\Quote{约匝达克的儿子}，他从未穿过污秽的衣服，总是穿着司铎的衣服，也从未被剥夺司铎职务。但那里所指的\Quote{耶叔亚}是基督，至高大父天主的司铎；他在\OriginalTerm{第一次来临}{first advent}时以谦卑、人性形体和可受难的方式而来，直到他的苦难；他自己也同样通过所有（苦难阶段）为我们众人成为祭品；他在复活后\BibleQuote{穿上长袍}，并被命名为天父的司铎直到永远。同样，我将对斋戒日习惯奉献的两只山羊做出解释。它们不也指向已经来临的基督特性的每个连续阶段吗？一对，一方面，相似（它们），因为主的普遍外貌的同一性，因为他不会以另一种形相来临，他必须被那些曾经伤害过他的人认出。但是其中一只，披着朱红布，在咒骂、普遍的唾弃、撕裂和刺穿中，被百姓丢弃在城外毁灭，带着基督受难的明显标记；他披着朱红衣服，被普遍唾弃，受尽一切侮辱，被钉在城外。另一只，为赎罪而献，作为食物\Emph{只给圣殿的司铎}，显明了第二次显现的明显证据；因为在所有罪过得到赎清之后，属灵圣殿的司铎，即教会的司铎，要享受主的恩宠的属灵公开分施（如同），而其他人都禁食不得救恩。\par

因此，由于关于\OriginalTerm{第一次来临}{first advent}的预言以多种形象遮蔽了它，并以各种羞辱贬低了它，而第二次（来临）则被预言为是显露的、全然相称于天主的，结果便是：他们只注视那一个他们容易理解和相信的（即第二次，在尊荣和荣耀中的），却（并非无根据地）被那更为隐晦——无论如何，更为卑微的——即第一次所欺骗。因此直到如今，他们仍坚称他们的基督尚未来临，因为祂没有在威严中来临；而他们却不知道，祂首先要以谦卑的姿态来临。\par

至此，我们已足够顺流而下，追溯基督行进的秩序，以此证明祂正如惯常所宣告的那样；为的是藉着圣经的这种和谐，我们能够明白；并且那些预言说将在基督之后发生的事件，可以相信已因天主的安排而成就。因为除非祂——在那之后这些事件必须成就——来临，否则那些预告其未来发生与祂来临相关的事件，绝不会发生。因此，如果你看见万民从此从人类错误的深渊中涌现，归向造物主天主和祂的基督（你不敢断言这没有被预言过，因为即使你如此断言，立刻——正如我们已经预先指出的——圣父的许诺会浮现在你面前，祂说：\BibleQuote{祢是我的儿子；我今日生祢。祢向我求，我就将外邦人赐祢\Emph{为}基业，将地极赐祢\Emph{为}田产。}\BibleRef{咏 2:7-8}而且你无法辩称那预言的对象是达味的儿子撒罗满，而非天主子基督；也无法辩称\Quote{地极}是应许给达味的儿子——他只统治犹大一地——而非天主子基督——祂已经用福音的光芒照亮了全世界。简言之，再者，一个\Quote{直到永远}的宝座更适合天主子基督，而非撒罗满——一个暂时的君王，只统治以色列。因为如今万民在呼求基督，他们过去并不认识祂；如今万民成群结队地投奔基督，他们过去对祂一无所知），你不能争辩说那是未来的事，而你却看见它在眼前发生。要么否认这些事件曾被预言，而它们就在你眼前；要么承认它们已被应验，而你亲耳听到它们被诵读；或者另一方面，如果你无法否认任一立场，那么它们便在那位关于祂而被预言的那位身上得着应验。\par

\SourceRecord{Translated by S. Thelwall.；Ante-Nicene Fathers；Vol. 3.；Edited by Alexander Roberts, James Donaldson, and A. Cleveland Coxe.；Buffalo, NY: Christian Literature Publishing Co.,；1885.；<http://www.newadvent.org/fathers/0308.htm>.}
